\documentclass{article}
\usepackage{amsmath,amssymb,amsthm}

\newtheorem{theorem}{Theorem}[subsection]
\newtheorem{lemma}[theorem]{Lemma}
\newtheorem*{mainthm}{Main Theorem}
\newtheorem{fact}{Fact}
\theoremstyle{definition}
\newtheorem{definition}{Definition}[section]
\numberwithin{equation}{section}

\begin{document}

\section{Recurrences}

\subsection{Fibonacci numbers}

\begin{definition}
The Fibonacci numbers are $F_0 = 0$, $F_1 = 1$ and
\begin{equation}\label{eq:fib}
  F_{n+2} = F_{n+1} + F_n \qquad (n \ge 0).
\end{equation}
\end{definition}

\begin{mainthm}
For every $n \ge 1$, $\gcd(F_n, F_{n+1}) = 1$.
\end{mainthm}

\begin{lemma}\label{lem:cassini}
For every $n \ge 1$,
\begin{equation}\label{eq:cassini}
  F_{n+1}F_{n-1} - F_n^2 = (-1)^n.
\end{equation}
\end{lemma}

\begin{proof}
We induct on $n$. For $n = 1$ the left side is $F_2F_0 - F_1^2 = -1$.
Assume \eqref{eq:cassini} for $n$. Using \eqref{eq:fib} twice,
\begin{align}
  F_{n+2}F_n - F_{n+1}^2
    &= (F_{n+1} + F_n)F_n - F_{n+1}(F_n + F_{n-1}) \notag\\
    &= F_n^2 - F_{n+1}F_{n-1} \notag\\
    &= -(-1)^n = (-1)^{n+1}. \qedhere
\end{align}
\end{proof}

\begin{proof}[Proof of the Main Theorem]
A common divisor of $F_n$ and $F_{n+1}$ divides the left side of
\eqref{eq:cassini}, hence divides $\pm 1$.
\end{proof}

\subsection{Strong induction}

\begin{theorem}\label{thm:zeck}
Every positive integer is a sum of distinct, non-consecutive Fibonacci
numbers.
\end{theorem}

\begin{proof}
Use strong induction on $n$. Let $F_k$ be the largest Fibonacci number with
$F_k \le n$. If $F_k = n$ we are done. Otherwise $0 < n - F_k < F_{k-1}$,
and by the induction hypothesis $n - F_k$ has such a representation, which
cannot use $F_{k-1}$.
\end{proof}

\begin{fact}
The ratio $F_{n+1}/F_n$ tends to the golden ratio
$\varphi = (1 + \sqrt5)/2$.
\end{fact}

\section{Binomial coefficients}

\begin{definition}
For integers $0 \le k \le n$, let
$\binom{n}{k} = \dfrac{n!}{k!\,(n-k)!}$.
\end{definition}

\begin{theorem}[Pascal's rule]
For $1 \le k \le n$,
\begin{equation}
  \binom{n+1}{k} = \binom{n}{k} + \binom{n}{k-1}.
\end{equation}
\end{theorem}

\begin{proof}
Count $k$-subsets of $\{1, \dots, n+1\}$ according to whether they contain
the element $n+1$.
\end{proof}

\begin{fact}
Summing Pascal's rule gives $\sum_{k=0}^{n} \binom{n}{k} = 2^n$.
\end{fact}

\end{document}
